% ============================================================
\chapter{Conclusion}
% ============================================================

\section{Summary of Achievements}

NetMind set out to let someone responsible for a small or medium branch
network describe what they want in natural language, in either Arabic or
English, rather than learn a vendor-specific command-line syntax. By the end
of Project 1, both of the system's operating modes were working and
demonstrated: a simulated mode that queries and controls a real Containerlab
network of routers, firewalls, hosts, and cameras, and a read-only
real-network mode that discovers devices on the physical local network by
combining four discovery methods, without altering them. Destructive actions
are guarded by a confirmation step, protected devices need an extra
acknowledgement, and every outcome is written to an audit log. A live
dashboard, an alert history with Telegram notifications, persistent storage,
and an automated test suite running on every change complete the system. In
a pilot evaluation of 28 English and Arabic requests (Section~\ref{sec:pilot}),
the model behaved correctly in 27 cases, answered every request in the
language it was asked in, and no destructive action ran without the user's
confirmation.

Because development started early, the team also completed the first three
milestones of the supervisor's improvement plan within this term:
stabilisation (October), the architecture refactor (November) and, ahead of
its December target, observability, which added a background collector and
diagnostic and alert-history tools for the agent.
Figure~\ref{fig:scope} places what was actually built against what remains
for Project 2, and Table~\ref{tab:coverage} summarises, capability area by
capability area, how the requirements gathered in Chapter~3 were met.

\begin{figure}[H]
\centering
\begin{tikzpicture}[font=\small,>={Stealth[length=2.5mm]}]
  \node[draw,rounded corners=3pt,fill=boxgray,text width=6.6cm,inner sep=7pt,align=left] (p1) {%
    {\large\bfseries\color{donegreen}Project 1 (achieved)}\par\smallskip
    \begin{itemize}[leftmargin=1.1em,itemsep=1pt,topsep=2pt]
      \item Simulated-network device status \& power control
      \item Firewall auditing \& port control (INPUT and FORWARD)
      \item Camera stream \& credential checks
      \item Real-network discovery: ARP, ping, IPv6, mDNS (read-only)
      \item Bilingual (EN/AR) chat with evidence view
      \item Live dashboard, alert history, Telegram alerts
      \item Background collector with UP/DEGRADED/DOWN states
      \item Diagnostic \& alert-history tools (21 tools)
      \item Confirmation step, protected devices \& audit log
      \item Services, backends, tool registry, model provider
      \item SQLite storage; server runs without root
      \item Automated tests, coverage \& continuous integration
    \end{itemize}};
  \node[draw,rounded corners=3pt,text width=6.0cm,inner sep=7pt,align=left,right=1.4cm of p1.north east,anchor=north west] (p2) {%
    {\large\bfseries\color{nextred}Project 2 (planned)}\par\smallskip
    \begin{itemize}[leftmargin=1.1em,itemsep=1pt,topsep=2pt]
      \item Fault injection \& root-cause analysis
      \item Proactive fixes, firewall rollback
      \item Authentication and roles
      \item MCP server and a hardened agent harness
      \item Benchmark, model comparison, usability study
    \end{itemize}};
  \draw[->,thick] (p1.east |- p2.west) -- (p2.west);
\end{tikzpicture}
\caption{Project 1 (achieved) scope against Project 2 (planned) scope}
\label{fig:scope}
\end{figure}

{\small
\begin{xltabular}{\textwidth}{|Y|L{3.4cm}|}
\hline
\head \textbf{Capability Area} & \textbf{Status} \\ \hline
\endhead
Device status \& power control (FR1, FR2) & Met \\ \hline
Network/segment visibility (FR3) & Met \\ \hline
Firewall auditing \& port control (FR4, FR5) & Met \\ \hline
Camera stream \& credential checks (FR6) & Met \\ \hline
Real-network discovery, read-only (FR7, NFR4) & Met \\ \hline
Bilingual Arabic/English interface (FR8) & Met \\ \hline
Scope guarding: refusing unrelated requests (FR9) & Met \\ \hline
Live monitoring dashboard (FR10) & Met \\ \hline
Confirmation before destructive actions, protected devices, audit log (FR11--FR13, NFR6) & Met \\ \hline
Connectivity check between lab devices (FR14) & Met \\ \hline
Alert notifications outside the browser (FR15) & Met \\ \hline
Usability, performance, portability (NFR1--NFR3, NFR9--NFR13) & Met \\ \hline
Least privilege, persistence, automated checks (NFR8, NFR14, NFR15) & Met \\ \hline
Agent behaviour in the pilot evaluation (Table~\ref{tab:pilot}) & 27/28 correct; 0 unconfirmed actions \\ \hline
Authentication (NFR7) & Planned for Project 2 \\ \hline
\caption{Requirement coverage summary}\label{tab:coverage}\\
\end{xltabular}}

\section{Limitations}

Three limitations were deliberately accepted in Project 1 rather than left
unnoticed. The first is that the web interface does not yet authenticate its
users (NFR7). The server only listens on the local machine and refuses
requests from other origins, which keeps it safe on a single laptop, but it
assumes one trusted operator.

The second is an architectural trade-off described in Chapter~3: device state
such as power status, firewall rules, and camera credentials is read live from
Docker rather than stored in a database of NetMind's own. Only very short
caches are used (three seconds for the device list, ten seconds for audits),
and they are cleared whenever a power state, rule or password changes, so
NetMind's picture cannot drift from the lab. This means every query needs a
reachable Docker daemon and live containers, which is acceptable for a
lab-scale deployment but would need to be revisited if NetMind were ever
asked to manage a larger or longer-lived network.

The third is that the agent's accuracy has only been measured on a small
scale. The pilot evaluation in Section~\ref{sec:pilot} used 28 prompts, one
run and one model; it gives encouraging results but cannot show how
consistent the model is across runs or how it compares with others. The same
pilot showed a median response time of about 9.5~seconds, which is longer
than intended and not yet broken down into model and tool time.

Two limitations listed earlier have since been resolved: conversation
history, once held only in memory and lost on restart, is now kept in the
SQLite storage described in Chapter~5, and alerts, once computed only while
the dashboard was open, are now raised by the background collector.

{\small
\begin{xltabular}{\textwidth}{|Y|Y|}
\hline
\head \textbf{Limitation} & \textbf{Planned Resolution} \\ \hline
\endhead
No authentication on the web interface & Authentication with viewer, operator and admin roles (\#31) \\ \hline
Device state read live from Docker & Kept for the lab; storing metrics over time is a possible extension \\ \hline
Accuracy measured only in a 28-prompt pilot; slow replies & Prompt benchmark with repeated runs, model comparison, latency breakdown and usability study (\#34--\#38) \\ \hline
\caption{Limitations and their planned resolution}\label{tab:limitations}\\
\end{xltabular}}

\section{Future Work}

Project 2 continues the supervisor's seven-month plan, which is tracked as
GitHub issues (issue numbers are given in brackets). Milestones M1 to M3 were
completed this term; the remaining milestones, adjusted to the academic
calendar, are listed in Table~\ref{tab:roadmap}.

{\small
\begin{xltabular}{\textwidth}{|L{3.4cm}|Y|}
\hline
\head \textbf{Milestone} & \textbf{Planned Work} \\ \hline
\endhead
M4: Diagnosis & Reversible fault-injection scenarios such as a link going down, packet loss, a stopped service or a firewall misconfiguration (\#25); root-cause-analysis playbooks that diagnose a fault step by step with evidence (\#26) \\ \hline
M5: Safe change & Proactive suggestions and remediation playbooks, always through the confirmation flow (\#29); firewall snapshots, diff, dry run and rollback (\#30); authentication and roles (\#31); new-device detection on the real network (\#32) \\ \hline
M6: Evaluation & A benchmark of more than 100 English and Arabic prompts measuring tool-selection accuracy, latency and safety (\#34, \#38); comparison of LLM providers and models (\#35); failure and load testing (\#36); a usability study (\#37) \\ \hline
M7: Delivery & Final thesis chapters, a v1.0 release, demo and mock defences (\#39--\#41) \\ \hline
\caption{Project 2 milestones}\label{tab:roadmap}\\
\end{xltabular}}

In addition to the plan, two pieces of work suggested by the supervisor
build directly on the structure delivered this term.

\textbf{An MCP server.} The Model Context Protocol (MCP) is an open standard
through which AI assistants connect to external tools. NetMind's tools
already exist as registered functions with generated schemas, capabilities
and risk levels, so they can be published through an MCP server. Any
MCP-compatible assistant could then query and manage the lab through the same
safety guard, confirmation flow and audit log, rather than only through
NetMind's own chat.

\textbf{A hardened agent harness.} The harness is the program around the
model that runs the loop: it sends the request, lets the model call a tool,
checks and runs the call, feeds the result back, and decides when to stop.
NetMind's agent loop is already a small harness with a round limit, a retry
safety net and a guard. Project 2 will strengthen it with clearer limits and
timeouts, structured logging of every step, and the prompt benchmark above,
so that changes to the model, the prompt or the tools can be measured rather
than judged by eye.

Figure~\ref{fig:timeline} places Project 1's phases on a timeline, for
reference against the Project 2 work that follows it.

\begin{figure}[H]
\centering
\begin{tikzpicture}[font=\footnotesize,x=3.3cm,y=0.85cm]
  \foreach \m/\i in {September/0,October/1,November/2,December/3}{
    \fill[headgray] (\i,0) rectangle (\i+1,-0.8);
    \node[font=\small\bfseries] at (\i+0.5,-0.4) {\m};
    \draw[black!25] (\i,0) -- (\i,7.0);
  }
  \draw[black!25] (4,0) -- (4,7.0);
  % \ph{start}{end}{row}{colour}{text}  (x in months from 1 September)
  % \ph{start}{end}{row}{colour}{text}: bar in months from 1 September, label to its right
  \newcommand{\ph}[5]{\draw[rounded corners=2pt,fill=#4] (#1,#3-0.3) rectangle (#2,#3+0.3);
    \node[anchor=west,font=\footnotesize] at (#2,#3) {#5};}
  \ph{0.0}{0.45}{0.8}{green!35}{Requirements \& design}
  \ph{0.25}{0.95}{1.7}{blue!25}{Simulated lab, firewalls, cameras}
  \ph{0.55}{1.25}{2.6}{orange!35}{Real mode \& live dashboard}
  \ph{1.0}{2.0}{3.5}{yellow!45}{M1: safety, tests, CI, \code{make up}}
  \ph{2.0}{2.75}{4.4}{red!30}{M2: refactor, SQLite, merged scan}
  \draw[rounded corners=2pt,fill=violet!25] (2.75,5.0) rectangle (2.98,5.6);
  \node[anchor=east,font=\footnotesize] at (2.75,5.3) {M3: background collector, diagnostic tools};
  \draw[rounded corners=2pt,fill=black!15] (2.7,5.9) rectangle (3.6,6.5);
  \node[anchor=east,font=\footnotesize] at (2.7,6.2) {Report \& presentation};
\end{tikzpicture}
\caption{Project 1 timeline, by phase}
\label{fig:timeline}
\end{figure}

Taken together, Project 1 delivered a working, demonstrated system covering
both of NetMind's operating modes, together with the safety, storage and
testing foundations that the planned work depends on. The remaining work is
to make NetMind diagnose faults, change the network safely, and measure how
well it does so; it extends what already exists
rather than redesigning it.
